\chapter{Introduction}
\label{chap:introduction}

\section{Project Context and Data Source}
Customer purchasing behaviour is an important input to retail decisions, yet transaction tables alone do not readily
show which customers buy repeatedly, which customers generate high value, or which customer groups require different
marketing approaches. CompanyX is the retail dataset used in this project. Its transactional schema records customers,
sales orders, order details, products, and associated customer demographics. The available sales history begins in May
2011 and ends in June 2014.

The project deliberately focuses on \emph{customer-centric analysis}, rather than general sales reporting. In
particular, the team uses Recency, Frequency, and Monetary (RFM) behaviour to describe customers, K-Means to discover
five interpretable groups, and a dimensional data warehouse to make these results available for analysis. The warehouse
also stores slowly changing customer attributes and monthly historical observations. Its automated ETL pipeline makes
the latest customer profiles and purchase-triggered segmentation results available at one-minute intervals.

\section{Business Problem}
The practical question is not simply how much CompanyX sold, but \emph{which customers contributed that value and how
  their purchasing patterns changed}. Four issues motivated the solution:
\begin{itemize}
  \item \textbf{Repeat purchasing and retention.} Single-order buyers and customers who do not return represent different opportunities from loyal repeat customers. An appropriate intervention requires a view of their purchase frequency and time since their last order.
  \item \textbf{Uneven customer value.} Customers differ considerably in their spending. Treating all customers identically can hide high-value accounts and lead to inefficient campaign allocation.
  \item \textbf{Cross-selling opportunities.} Understanding a customer's value and engagement stage can help marketers decide when to recommend accessories, maintenance products, or complementary items. RFM identifies target groups; verifying product-specific cross-selling gaps would require separate product-level analysis.
  \item \textbf{Fragmented historical and operational views.} A record updated when a purchase occurs answers a different question from a month-end view of the entire customer population. Both are needed to distinguish purchase-time behaviour from longer-term inactivity and segment migration.
\end{itemize}

The Power BI dashboard provides key indicators such as repeat-customer rate, customer inactivity, average purchase
frequency, and spending patterns. Together, these metrics offer an overview of customer purchasing behavior and
engagement, helping decision-makers identify potential retention issues, understand customer value, and develop
appropriate marketing strategies.

\section{Objectives and Deliverables}

The primary objective of this project is to develop a customer segmentation and business intelligence system that helps
CompanyX better understand customer purchasing behaviour and support data-driven decision-making. The project focuses
on two main goals:

\begin{enumerate}
  \item \textbf{Customer Segmentation:}
        Analyse customer purchasing patterns to identify distinct customer
        groups based on their behaviour and value. These insights can help
        the company better understand its customer base and develop more
        effective marketing, retention, and customer engagement strategies.

  \item \textbf{Automated Business Intelligence Pipeline:}
        Develop an integrated system that automatically collects, processes,
        and analyses customer data. The system aims to provide up-to-date
        customer information while maintaining historical records, allowing
        businesses to monitor changes in customer behaviour over time.
\end{enumerate}

To achieve these objectives, the project delivers an integrated solution consisting of a data warehouse, automated data
processing pipelines, machine learning-based customer segmentation, and interactive business intelligence dashboards.
Together, these components transform raw transactional data into meaningful insights that support business analysis and
decision-making.
